% Image credits for the photographs used in this volume (Book 2 only).
% Machine-readable license records live in images/book2/CREDITS.md (and
% images/book1/CREDITS.md for the three photographs reused from Book 1);
% keep the files in sync when adding or removing a photograph.
\chapter*{Kredit Gambar}

% Under bidi=basic a Latin run is ONE directional box, so a long author or
% institution name cannot be broken at all; \allowbreak inside it does nothing.
% The Arabic edition reported this page as the book's only overfull boxes.
% \sloppy must be in force when the paragraph is BROKEN, hence the \par before
% \endgroup -- without it the group closes first and the tolerance is restored.
\begingroup\sloppy

Gambar-gambar dalam buku ini adalah karya asli. Foto-fotonya digunakan
berdasarkan lisensi bebas masing-masing, dengan terima kasih kepada
para pembuatnya:

\begin{itemize}
  \item \emph{Gregor Mendel} (bab meiosis): foto, abad ke-19,
        Wikimedia Commons, domain publik.
  \item \emph{Alexander Fleming} (bab resistensi antibiotik): foto,
        Imperial War Museums, 1943, melalui Wikimedia Commons, domain
        publik.
  \item \emph{HIV yang bertunas dari limfosit} (bab kekebalan
        adaptif): mikrograf elektron karya C.~Goldsmith, Centers for
        Disease Control and Prevention, melalui Wikimedia Commons,
        domain publik.
  \item \emph{Burung finch Darwin} (bab diversifikasi): ukiran karya
        John Gould, 1845, Wikimedia Commons, domain publik.
  \item \emph{Homunkulus sensorik} (bab gerak): ilustrasi dari
        OpenStax \emph{Anatomy and Physiology}, melalui Wikimedia
        Commons, CC~BY~3.0.
  \item \emph{Charles Darwin} (bab pohon filogenetik): foto karya
        Henry Maull dan John Fox, sekitar 1854, Wikimedia Commons,
        domain publik.
  \item \emph{Edward Jenner} (bab kekebalan adaptif): lukisan cat
        minyak, Wellcome Collection melalui Wikimedia Commons,
        CC~BY~4.0.
  \item \emph{Fosil Archaeopteryx} (bab nenek moyang bersama): foto
        karya James L.~Amos, Wikimedia Commons, CC0.
\end{itemize}

Tautan sumber lengkap dan alamat lisensi setiap foto tercantum di
\texttt{images/\allowbreak book2/\allowbreak CREDITS.md} dan \texttt{images/\allowbreak book1/\allowbreak CREDITS.md}
di dalam repositori buku ini.

\bigskip
Selain foto dan gambar asli tersebut, sekitar lima puluh gambar yang
tersebar di seluruh bab jilid ini adalah ilustrasi hasil kecerdasan
buatan, dibuat dengan pembangkit gambar OpenAI dari prompt yang ditulis
oleh para penyunting buku, dan diperiksa ketepatan biologisnya sebelum
disertakan. Prompt lengkap setiap gambar tercatat dalam
\texttt{images/\allowbreak book2/\allowbreak ai/\allowbreak PROMPTS.md} di repositori.
\par\endgroup
\clearpage
